% Generated by publication/build_sections.py; do not edit by hand.
\paragraph{Informal verdict.}
\textbf{A --- Complete solution}, with \textbf{high confidence}.

\paragraph{Lean coverage.}
For every strictly positive summable real sequence, Lean defines the full achievement set of all Boolean-selected subseries, including finite selections, and proves the exact finite-stage characterization: Sigma\(a\) is the intersection of F\_n + [0,r\_n]; r\_n + Delta\_n is the attained longest interval length in the n-th approximant and decreases to the attained longest interval length in Sigma\(a\); hence interior Sigma\(a\) is empty iff Delta\_n tends to zero, equivalently iff for every epsilon > 0 some r\_n + Delta\_n < epsilon. This is unconditional and stronger than the printed hypothesis because no monotonicity assumption is needed, so it specializes directly to every positive decreasing summable sequence in Q803.

\subsection{Audited informal answer}
Q803 --- Exact finite-stage characterization
\begin{enumerate}
\item Faithful restatement
\end{enumerate}
Let

\[
a_1\ge a_2\ge \cdots>0,
\qquad 
\sum_{n=1}^{\infty}a_n<\infty .
\]

Define the set of all subsums

\[
\Sigma(a):=
\left\{
\sum_{n=1}^{\infty}\varepsilon_n a_n:
\varepsilon_n\in\{0,1\}
\right\}.
\]

Finite extracted series correspond to eventually-zero sequences \((\varepsilon_n)\). Including or excluding the empty subseries, whose sum is \(0\), does not affect the interior.
The question asks for a characterization of those sequences for which

\[
\operatorname{int}\Sigma(a)=\varnothing.
\]

This is exactly the complete Q803 statement in the PDF; there are no subparts, footnotes, or further qualifications. In standard French usage, ``décroissante'' means nonincreasing. There is no apparent typo or substantive ambiguity. Luc Pommeret

\begin{enumerate}
\item The exact characterization
\end{enumerate}
For \(n\ge 0\), put

\[
r_n:=\sum_{k>n}a_k,
\qquad r_0=\sum_{k\ge1}a_k,
\]

and let

\[
F_n:=
\left\{
\sum_{k=1}^{n}\varepsilon_k a_k:
\varepsilon_k\in\{0,1\}
\right\}
\]

be the finite set of subsums of the first \(n\) terms. Write its distinct elements in increasing order:

\[
F_n=\{f_{n,0}<f_{n,1}<\cdots<f_{n,m_n}\}.
\]

Partition this ordered set into maximal consecutive blocks

\[
f_{n,p},f_{n,p+1},\ldots,f_{n,q}
\]

such that

\[
f_{n,j+1}-f_{n,j}\le r_n
\qquad (p\le j<q).
\]

Thus a new block begins whenever a consecutive gap is strictly larger than \(r_n\).
Define

\[
\Delta_n:=
\max_{\text{such maximal blocks }[p,q]}
\bigl(f_{n,q}-f_{n,p}\bigr).
\]

Theorem
The sequence

\[
\ell_n:=r_n+\Delta_n
\]

is nonincreasing and converges. Moreover,

\[
\boxed{
\lim_{n\to\infty}\ell_n
=
\max\bigl\{\,b-a:[a,b]\subseteq\Sigma(a)\,\bigr\}.
}
\]

In particular, because \(r_n\to0\),

\[
\boxed{
\operatorname{int}\Sigma(a)=\varnothing
\quad\Longleftrightarrow\quad
\lim_{n\to\infty}\Delta_n=0.
}
\]

Equivalently,

\[
\boxed{
\operatorname{int}\Sigma(a)=\varnothing
\quad\Longleftrightarrow\quad
\forall\varepsilon>0\ \exists n\quad
r_n+\Delta_n<\varepsilon .
}
\]

Thus the answer is: the subsum set has empty interior precisely when every \(r_n\)-connected cluster of finite subsums of the first \(n\) terms has diameter tending to zero.
This finite-chain criterion is essentially the analytic characterization recorded as Proposition 2.8 in the achievement-set literature. arXiv+1

\begin{enumerate}
\item Proof of the characterization
\end{enumerate}
3.1. Compactness of the subsum set
Let

\[
\Omega:=\{0,1\}^{\mathbb N}
\]

and define

\[
\pi:\Omega\longrightarrow \mathbb R,
\qquad
\pi(\varepsilon):=\sum_{n=1}^{\infty}\varepsilon_n a_n.
\]

Equip \(\Omega\) with the usual product metric, for example

\[
d(\varepsilon,\eta)
=
\sum_{n=1}^{\infty}2^{-n}|\varepsilon_n-\eta_n|.
\]

If \(\varepsilon_k=\eta_k\) for \(1\le k\le N\), then

\[
|\pi(\varepsilon)-\pi(\eta)|
\le \sum_{k>N}a_k=r_N.
\]

Since \(r_N\to0\), the map \(\pi\) is uniformly continuous.
The space \(\Omega\) is compact: from any sequence in \(\Omega\), one first extracts a subsequence on which the first coordinate is constant, then a further subsequence on which the second coordinate is constant, and so on; the diagonal subsequence converges. Hence

\[
\Sigma(a)=\pi(\Omega)
\]

is compact.

3.2. Canonical finite approximations
Define

\[
I_n:=F_n+[0,r_n]
=
\bigcup_{f\in F_n}[f,f+r_n].
\]

Each \(I_n\) is a finite union of closed bounded intervals.
We claim that

\[
I_{n+1}\subseteq I_n.
\]

Indeed,

\[
F_{n+1}=F_n\cup(F_n+a_{n+1}),
\qquad
r_n=a_{n+1}+r_{n+1}.
\]

For each \(f\in F_n\), the two child intervals at level \(n+1\) are

\[
[f,f+r_{n+1}]
\]

and

\[
[f+a_{n+1},f+a_{n+1}+r_{n+1}]
=
[f+a_{n+1},f+r_n].
\]

Both are contained in \([f,f+r_n]\). Therefore \(I_{n+1}\subseteq I_n\).
We now prove the fundamental identity

\[
\boxed{\Sigma(a)=\bigcap_{n=0}^{\infty}I_n.}
\]

The inclusion \(\Sigma(a)\subseteq I_n\) is immediate: if

\[
x=\sum_{k=1}^{\infty}\varepsilon_k a_k,
\]

then

\[
x=
\underbrace{\sum_{k=1}^{n}\varepsilon_k a_k}_{\in F_n}
+
\underbrace{\sum_{k>n}\varepsilon_k a_k}_{\in[0,r_n]}.
\]

Conversely, suppose \(x\in I_n\) for every \(n\). For each \(n\), choose a binary word

\[
w^{(n)}=(\varepsilon^{(n)}_1,\ldots,\varepsilon^{(n)}_n)
\]

such that

\[
0\le
x-\sum_{k=1}^{n}\varepsilon_k^{(n)}a_k
\le r_n.
\]

Extend \(w^{(n)}\) by zeros to an element \(\widetilde w^{(n)}\in\Omega\). By the diagonal compactness argument, a subsequence

\[
\widetilde w^{(n_j)}
\]

converges to some \(\varepsilon\in\Omega\). Continuity of \(\pi\) gives

\[
\pi\bigl(\widetilde w^{(n_j)}\bigr)\longrightarrow \pi(\varepsilon).
\]

But

\[
\left\lvert{}
x-\pi\bigl(\widetilde w^{(n_j)}\bigr)
\right\rvert{}
\le r_{n_j}\longrightarrow0,
\]

so \(\pi(\varepsilon)=x\). Hence \(x\in\Sigma(a)\).

3.3. A lemma on decreasing finite unions of intervals
We use the following general fact.
Lemma
Let

\[
W_0\supseteq W_1\supseteq W_2\supseteq\cdots
\]

be compact subsets of \(\mathbb R\), each of which is a finite union of closed intervals. Let \(L_n\) be the maximum length of a connected component of \(W_n\). Then \(L_n\) is nonincreasing, and if

\[
W:=\bigcap_{n=0}^{\infty}W_n,
\qquad
L:=\lim_{n\to\infty}L_n,
\]

then

\[
L=\max\{b-a:[a,b]\subseteq W\}.
\]

Proof
Every component of \(W_{n+1}\), being connected and contained in \(W_n\), lies in a single component of \(W_n\). Therefore

\[
L_{n+1}\le L_n.
\]

Suppose first that \([a,b]\subseteq W\). For every \(n\), the connected interval \([a,b]\) lies in one component of \(W_n\). Hence

\[
b-a\le L_n
\]

for every \(n\), and therefore

\[
b-a\le L.
\]

It remains to construct an interval of length \(L\) in \(W\). For each \(n\), choose the leftmost component

\[
C_n=[u_n,v_n]
\]

of \(W_n\) whose length is \(L_n\). All pairs \((u_n,v_n)\) lie in a fixed compact rectangle. Passing to a subsequence,

\[
u_{n_j}\longrightarrow u,
\qquad
v_{n_j}\longrightarrow v.
\]

Since

\[
v_{n_j}-u_{n_j}=L_{n_j}\longrightarrow L,
\]

we have

\[
v-u=L.
\]

Fix \(m\). For sufficiently large \(j\), \(n_j\ge m\), so

\[
C_{n_j}\subseteq W_{n_j}\subseteq W_m.
\]

For \(t\in[0,1]\), the point

\[
x_j(t):=(1-t)u_{n_j}+tv_{n_j}
\]

belongs to \(W_m\), and

\[
x_j(t)\longrightarrow (1-t)u+tv.
\]

Since \(W_m\) is closed,

\[
(1-t)u+tv\in W_m.
\]

Thus \([u,v]\subseteq W_m\) for every \(m\), and hence

\[
[u,v]\subseteq W.
\]

Its length is \(L\). \(\square\)

3.4. Computing the components of \(I_n\)
Recall

\[
I_n=\bigcup_{j=0}^{m_n}[f_{n,j},f_{n,j}+r_n].
\]

Two consecutive intervals overlap or touch exactly when

\[
f_{n,j+1}\le f_{n,j}+r_n,
\]

that is,

\[
f_{n,j+1}-f_{n,j}\le r_n.
\]

Consequently, if

\[
f_{n,p},\ldots,f_{n,q}
\]

is one maximal \(r_n\)-connected block, the corresponding component of \(I_n\) is

\[
[f_{n,p},f_{n,q}+r_n],
\]

whose length is

\[
f_{n,q}-f_{n,p}+r_n.
\]

Therefore the maximum component length of \(I_n\) is exactly

\[
L_n=r_n+\Delta_n=\ell_n.
\]

Applying the preceding lemma to the decreasing sequence \((I_n)\), and using

\[
\Sigma(a)=\bigcap_n I_n,
\]

we obtain

\[
\max\{b-a:[a,b]\subseteq\Sigma(a)\}
=
\lim_{n\to\infty}(r_n+\Delta_n).
\]

Finally, \(r_n\to0\), so

\[
\lim_{n\to\infty}\Delta_n
=
\lim_{n\to\infty}(r_n+\Delta_n).
\]

The set \(\Sigma(a)\) has nonempty interior exactly when it contains a nondegenerate interval. This proves

\[
\operatorname{int}\Sigma(a)=\varnothing
\iff
\Delta_n\longrightarrow0.
\]

\(\square\)

\begin{enumerate}
\item Consequence: the empty-interior case is exactly the Cantor-set case
\end{enumerate}
Under the original hypotheses, \(\Sigma(a)\) has no isolated points.
Indeed, take

\[
x=\sum_{n=1}^{\infty}\varepsilon_n a_n\in\Sigma(a).
\]

Either infinitely many \(\varepsilon_n\) are \(0\), in which case one can change a sufficiently far \(0\) into a \(1\), or infinitely many are \(1\), in which case one can change a sufficiently far \(1\) into a \(0\). This produces a distinct element of \(\Sigma(a)\) at distance \(a_n\to0\) from \(x\).
Thus \(\Sigma(a)\) is compact and perfect. If its interior is empty, it is nowhere dense. Moreover, any connected subset of \(\mathbb R\) containing two distinct points contains the whole interval between them; hence an empty-interior \(\Sigma(a)\) is totally disconnected.
Therefore, under the hypotheses of Q803,

\[
\boxed{
\operatorname{int}\Sigma(a)=\varnothing
\quad\Longleftrightarrow\quad
\Sigma(a)\text{ is a Cantor set}.
}
\]


\begin{enumerate}
\item The classical term-versus-tail regimes
\end{enumerate}
The exact criterion above is needed because comparisons between \(a_n\) and \(r_n\) do not decide the mixed case. They do, however, settle the two extreme regimes.
5.1. If the tail eventually dominates the term
Suppose

\[
a_n\le r_n
\]

for all \(n>N\).
At such an index \(n\), for each \(f\in F_{n-1}\),

\[
[f,f+r_n]
\quad\text{and}\quad
[f+a_n,f+a_n+r_n]
\]

overlap, and their union is

\[
[f,f+a_n+r_n]=[f,f+r_{n-1}].
\]

Hence

\[
I_n=I_{n-1}.
\]

It follows that

\[
I_n=I_N\qquad(n\ge N),
\]

and therefore

\[
\Sigma(a)=I_N.
\]

Thus \(\Sigma(a)\) is a finite union of nondegenerate closed intervals and has nonempty interior.
In particular,

\[
\boxed{
\operatorname{int}\Sigma(a)=\varnothing
\ \Longrightarrow\\
a_n>r_n\text{ for infinitely many }n.
}
\]

If \(a_n\le r_n\) for every \(n\), then \(I_n=I_0\) for every \(n\), so

\[
\boxed{
\Sigma(a)=\left[0,\sum_{n=1}^{\infty}a_n\right].
}
\]


5.2. If the term eventually dominates the tail
Suppose

\[
a_n>r_n
\]

for all \(n>N\).
Consider the tail achievement set

\[
\Sigma_N:=
\left\{
\sum_{n>N}\varepsilon_n a_n:\varepsilon_n\in\{0,1\}
\right\}.
\]

Its natural binary coding is injective. Indeed, if two binary sequences first differ at index \(j>N\), say \(\varepsilon_j=1\) and \(\eta_j=0\), then

\[
\sum_{n>N}(\varepsilon_n-\eta_n)a_n
\ge a_j-\sum_{n>j}a_n
=a_j-r_j>0.
\]

Thus \(\Sigma_N\) is homeomorphic to the binary Cantor space and in particular contains no interval.
Now

\[
\Sigma(a)=F_N+\Sigma_N
\]

is a finite union of translates of the compact nowhere-dense set \(\Sigma_N\). Hence it too has empty interior.
Therefore

\[
\boxed{
a_n>r_n\text{ eventually}
\ \Longrightarrow\\
\operatorname{int}\Sigma(a)=\varnothing.
}
\]

The converse is false: empty interior can occur even though \(a_n\le r_n\) infinitely often.

\begin{enumerate}
\item Why the mixed case cannot be decided by the signs of \(a_n-r_n\)
\end{enumerate}
Consider, for \(0<q\le 2/3\), the strictly decreasing multigeometric sequence

\[
a_{2k+1}=3q^k,
\qquad
a_{2k+2}=2q^k,
\qquad k\ge0.
\]

The tail after the first term in the \(k\)-th block is

\[
r_{2k+1}
=
2q^k+\sum_{j>k}5q^j
=
q^k\left(2+\frac{5q}{1-q}\right),
\]

whereas the tail after the second term is

\[
r_{2k+2}
=
\sum_{j>k}5q^j
=
q^k\frac{5q}{1-q}.
\]

For both \(q=\frac15\) and \(q=\frac14\), one has

\[
a_{2k+1}<r_{2k+1},
\qquad
a_{2k+2}>r_{2k+2}.
\]

Thus the complete pattern of inequalities is identical: every odd index is tail-dominated, and every even index is term-dominated.
Nevertheless, the corresponding achievement sets have opposite interior behavior.

6.1. The case \(q=\frac15\): empty interior
After \(m\) complete blocks there are at most \(4^m\) distinct prefix subsums, because each block contributes one of

\[
0,\ 2,\ 3,\ 5
\]

times the appropriate power of \(1/5\).
The remaining tail is

\[
r_{2m}
=
\sum_{k=m}^{\infty}5\cdot5^{-k}
=
\frac{25}{4}\,5^{-m}.
\]

Thus \(I_{2m}\) is covered by at most \(4^m\) intervals, each of length \(r_{2m}\). Its largest component therefore has length at most the sum of these lengths:

\[
\ell_{2m}
\le
4^m r_{2m}
=
\frac{25}{4}\left(\frac45\right)^m
\longrightarrow0.
\]

Since \((\ell_n)\) is nonincreasing,

\[
\ell_n\longrightarrow0.
\]

By the theorem,

\[
\boxed{
\operatorname{int}\Sigma\!\left(3,2;\frac15\right)=\varnothing.
}
\]


6.2. The case \(q=\frac14\): a nondegenerate interval
Now let

\[
a_{2k+1}=\frac{3}{4^k},
\qquad
a_{2k+2}=\frac{2}{4^k}.
\]

A complete block contributes a digit in

\[
D:=\{0,2,3,5\}.
\]

After \(m\) blocks, the finite block subsums are

\[
B_m
=
\left\{
\sum_{k=0}^{m-1}\frac{d_k}{4^k}:d_k\in D
\right\}.
\]

Scale them by \(4^{m-1}\):

\[
A_m:=4^{m-1}B_m
=
\left\{
\sum_{k=0}^{m-1}d_k4^{m-1-k}:d_k\in D
\right\}
\subseteq\mathbb Z.
\]

We claim that \(A_m\) contains every integer in the interval

\[
[\alpha_m,\beta_m],
\qquad
\alpha_m:=\frac{2(4^m-1)}3,
\qquad
\beta_m:=4^m-1.
\]

For \(m=1\),

\[
A_1=D
\]

contains \(2,3\), and indeed

\[
\alpha_1=2,\qquad\beta_1=3.
\]

Suppose the claim holds for \(m\). Since

\[
A_{m+1}=4A_m+D,
\]

the set \(A_{m+1}\) contains every integer from

\[
4\alpha_m+2
\]

to

\[
4\beta_m+3.
\]

To see this, consider residues modulo \(4\):


residues \(2\) and \(3\) are obtained using the digits \(2\) and \(3\);


a multiple \(4t\) in the stated range is obtained using digit \(0\);


a number \(4t+1\) is written as \(4(t-1)+5\).


The endpoint recurrences are

\[
4\alpha_m+2=\alpha_{m+1},
\qquad
4\beta_m+3=\beta_{m+1}.
\]

This proves the claim.
The tail after \(m\) blocks is

\[
r_{2m}
=
\sum_{k=m}^{\infty}\frac5{4^k}
=
\frac{20}{3}\,4^{-m}.
\]

After multiplying by \(4^{m-1}\), its length becomes

\[
4^{m-1}r_{2m}=\frac53.
\]

Since \(A_m\) contains every consecutive integer from \(\alpha_m\) through \(\beta_m\), the intervals of length \(5/3\) starting at these integers overlap and cover

\[
\left[\alpha_m,\beta_m+\frac53\right].
\]

Scaling back, \(I_{2m}\) contains

\[
J_m
=
\left[
\frac{\alpha_m}{4^{m-1}},
\frac{\beta_m+5/3}{4^{m-1}}
\right].
\]

The endpoints simplify to

\[
J_m
=
\left[
\frac83\left(1-4^{-m}\right),
\,
4+\frac83\,4^{-m}
\right].
\]

Hence

\[
\left[\frac83,4\right]\subseteq J_m\subseteq I_{2m}
\]

for every \(m\).
Because the even-indexed approximants are cofinal in the decreasing family \((I_n)\),

\[
\bigcap_{m=1}^{\infty}I_{2m}
=
\bigcap_{n=1}^{\infty}I_n
=
\Sigma(a).
\]

Therefore

\[
\boxed{
\left[\frac83,4\right]
\subseteq
\Sigma\!\left(3,2;\frac14\right).
}
\]

In particular,

\[
\boxed{
\operatorname{int}\Sigma\!\left(3,2;\frac14\right)\ne\varnothing.
}
\]

Thus even the complete pattern of inequalities

\[
a_n\lessgtr r_n
\]

does not characterize empty interior.

\begin{enumerate}
\item Boundary tests
\end{enumerate}
7.1. Geometric sequences
Let

\[
a_n=q^{n-1},
\qquad 0<q<1.
\]

Then

\[
r_n=\frac{q^n}{1-q}.
\]

If \(q<\frac12\), then

\[
a_n>r_n
\]

for every \(n\), so the achievement set has empty interior.
If \(q\ge\frac12\), then

\[
a_n\le r_n
\]

for every \(n\), and therefore

\[
\Sigma(a)
=
\left[0,\frac1{1-q}\right].
\]

At the boundary \(q=\frac12\), equality \(a_n=r_n\) holds at every index, and

\[
\Sigma(a)=[0,2].
\]

This shows that the strictness of the tail comparison matters.

7.2. Finite support
If only finitely many \(a_n\) are nonzero, then for sufficiently large \(n\),

\[
r_n=0,
\qquad
I_n=F_n.
\]

The achievement set is finite and therefore has empty interior. The original assumption \(a_n>0\) for every \(n\) excludes this degenerate case.

\begin{enumerate}
\item A useful easily checked sufficient condition
\end{enumerate}
Since \(I_n\) is a union of at most \(2^n\) intervals of length \(r_n\), its largest component satisfies

\[
\ell_n\le 2^nr_n.
\]

Consequently,

\[
\boxed{
\liminf_{n\to\infty}2^nr_n=0
\quad\Longrightarrow\quad
\operatorname{int}\Sigma(a)=\varnothing.
}
\]

More sharply, one may replace \(2^n\) by the actual number \(|F_n|\):

\[
\liminf_{n\to\infty}|F_n|\,r_n=0
\quad\Longrightarrow\quad
\operatorname{int}\Sigma(a)=\varnothing.
\]

These are sufficient, not necessary, conditions.

\begin{enumerate}
\item Exact finite-stage computation
\end{enumerate}
For rational input data and exactly known tails, the numbers

\[
\ell_n=r_n+\Delta_n
\]

can be computed without numerical tolerances.
PythonExécuterfrom fractions import Fraction
from collections.abc import Sequence, Iterator


def interval\_moduli(
    terms: Sequence[Fraction],
    tails: Sequence[Fraction],
) -> Iterator[tuple[int, Fraction]]:
    """
    Compute ell\_n = the largest component length of

        I\_n = F\_n + [0, r\_n],

    where:
        terms[n-1] = a\_n,
        tails[n]   = r\_n = sum\_\{k>n\} a\_k.

    The array tails must have length len\(terms\) + 1.
    Use Fraction values to make all comparisons exact.
    """
    terms = [Fraction\(x\) for x in terms]
    tails = [Fraction\(x\) for x in tails]

    if len\(tails\) != len\(terms\) + 1:
        raise ValueError("tails must have length len\(terms\) + 1")

    finite\_sums = \{Fraction\(0\)\}

    for n, a\_n in enumerate\(terms, start=1\):
        old\_sums = tuple\(finite_sums\)
        finite\_sums.update\(x + a_n for x in old_sums\)

        values = sorted\(finite_sums\)
        r\_n = tails[n]

        left = values[0]
        previous = values[0]
        best = Fraction\(0\)

        for x in values[1:]:
            if x - previous <= r\_n:
                previous = x
            else:
                best = max\(best, previous - left + r_n\)
                left = previous = x

        best = max\(best, previous - left + r_n\)
        yield n, best
At stage \(n\), the worst-case number of distinct finite subsums is \(2^n\), so the straightforward method is exponential. Finite computations can reveal candidate behavior, but a proof of empty interior still requires a rigorous estimate forcing \(\ell_n\to0\), while a proof of nonempty interior requires a persistent positive lower bound or an explicitly exhibited common interval.

\begin{enumerate}
\item Hypothesis and topology audit
\end{enumerate}
Topological assumptions
Everything takes place in \(\mathbb R\) with its usual Euclidean topology. It is Hausdorff, second-countable, complete, and metrizable. Connected subsets of \(\mathbb R\) are intervals; this is the only order-topological fact used about connectedness.
Compactness
Compactness of \(\{0,1\}^{\mathbb N}\) was used only in its countable metric form and was proved by a diagonal subsequence argument. Compactness of bounded endpoint sequences was used through Bolzano--Weierstrass. No nets, filters, or general uncountable Tychonoff theorem are needed.
Countability
The proof is entirely sequential and depends on the sequence being indexed by \(\mathbb N\). There are no uncountable families of subsums or covers.
Choice
Within ordinary ZFC, no substantial form of the axiom of choice is used. Choices from finite sets may be made canonically---for example, by taking the lexicographically least coding word or the leftmost longest component. No uncountable choice is involved.
Baire category and measure
The main characterization does not use the Baire category theorem or Lebesgue measure. The finite-union argument in the eventually fast regime uses only the elementary fact that a finite union of closed nowhere-dense sets is nowhere dense.
Necessity of the stated sequence hypotheses


Summability is essential to this formulation. It gives \(r_n\to0\), convergence of every subseries, uniform convergence of the coding map, and compactness of \(\Sigma(a)\). For a divergent positive series, the set of finite subsums is unbounded and infinite subseries may diverge.


Monotonicity is not needed for the exact criterion. The proof works for every summable nonnegative sequence in its given order. Since the subsum set is invariant under permutations, monotonicity merely gives a canonical ordering.


Strict positivity can be weakened to nonnegativity. Zero terms may simply be deleted. Strict positivity is used when asserting that an infinite achievement set has no isolated points.


Even signed absolutely summable sequences reduce to the positive case. If \(\sum |a_n|<\infty\), then

\[
\Sigma(a)
=
\sum_{a_n<0}a_n+\Sigma(|a_n|),
\]

because at a negative coordinate selecting \(a_n\) is equivalent, after translation, to not selecting \(|a_n|\). Thus signed achievement sets are translates of positive ones and have the same interior behavior.



\begin{enumerate}
\item Final classification
\end{enumerate}
Define \(r_n\), \(F_n\), and \(\Delta_n\) as above. Then

\[
\boxed{
\operatorname{int}\Sigma(a)=\varnothing
\iff
\lim_{n\to\infty}\Delta_n=0,
}
\]

where \(\Delta_n\) is the largest span of a maximal chain of consecutive \(n\)-term subsums whose successive gaps are at most the remaining tail \(r_n\).
More strongly,

\[
\boxed{
\max\{\,|J|:J\subseteq\Sigma(a)\text{ is an interval}\,\}
=
\lim_{n\to\infty}\bigl(r_n+\Delta_n\bigr).
}
\]

This is a complete literal answer to Q803. It is, however, an exponentially complicated finite-stage criterion rather than a simple condition involving only the individual comparisons \(a_n\) versus \(r_n\). Current surveys distinguish this exact analytic characterization from the still-unresolved search for an easy structural test separating Cantor sets from Cantorvals in the mixed regime. arXiv+1
